\section*{Chapter \ref{ch:m3:prediction-and-betting-markets} --- Prediction and Betting Markets}

\begin{solution}{exo:m3:prediction-and-betting-markets:1}
$1/1.80 + 1/2.10 = 1.0317$: an \omterm{def:m3:prediction-and-betting-markets:book}{overround} of 3.17\%; multiplicative probabilities 53.85\% and 46.15\%.
\end{solution}

\begin{solution}{exo:m3:prediction-and-betting-markets:2}
Lay $100 \times 3.0/2.5 = 120$ at 2.5: if the selection wins, $+200 - 180 = 20$; if it loses, $+120 - 100 = 20$. A profit of 20 either way, before
commission.
\end{solution}

\begin{solution}{exo:m3:prediction-and-betting-markets:3}
\omterm{def:m3:prediction-and-betting-markets:event}{Event contracts} are derivatives under the Commodity Exchange Act, and the statute requires an entity listing them for public trading to be designated as
a contract market by the CFTC.
\end{solution}

\begin{solution}{exo:m3:prediction-and-betting-markets:4}
64.82\%, 22.71\% and 12.48\%. Raising each \omterm{def:m3:prediction-and-betting-markets:book}{implied probability} to a common power above one shrinks small probabilities proportionally more, so the
longshot loses more than under proportional scaling (13.48\%).
\end{solution}

\begin{solution}{exo:m3:prediction-and-betting-markets:5}
Decimal odds $1/0.58 = 1.724$: $(0.65 \times 1.724 - 1)/0.724 = 16.7\%$ of wealth.
\end{solution}

\begin{solution}{exo:m3:prediction-and-betting-markets:6}
Each venue resolves by its own rules and source; a disputed, delayed or ambiguous outcome can resolve YES on one and NO on the other, or be voided on one, so
the two positions no longer offset.
\end{solution}

\begin{solution}{exo:m3:prediction-and-betting-markets:7}
About $-3.1\%$ between 1 and 3, and $-42.1\%$ between 50 and 1\,000.
\end{solution}

\begin{solution}{exo:m3:prediction-and-betting-markets:8}
50 to 1 is decimal odds of 51, an \omterm{def:m3:prediction-and-betting-markets:book}{implied probability} of about 2\% before the margin, and the margin and the \omterm{def:m3:prediction-and-betting-markets:flb}{favourite--longshot bias} sit
disproportionately on longshots: the true probability is likely lower, perhaps much lower.
\end{solution}

\begin{solution}{pb:m3:prediction-and-betting-markets:1}
\textbf{1.} $0.58 + 0.38 + 0.01 + 0.005 = 0.975$. \textbf{2.} Exactly one dollar, at resolution, if both venues resolve the same way. \textbf{3.} $1 - 0.975
\times (1 + 4.5\% \times 30/365) = 0.0214$ a pair. \textbf{4.} About 1.86 cents: the two fees plus a month of capital. \textbf{5.} About USD~214, locking
USD~9\,750 for a month (and margin or collateral on each venue). \textbf{6.} The pair can pay nothing or two dollars, or one leg can be voided: the
arbitrage becomes a bet on the rules. \textbf{7.} \omterm{def:m3:blockchains-for-traders:stable}{Stablecoin} risk (\cref{ch:m3:blockchains-for-traders}) and the cost of moving dollars on and off chain.
\textbf{8.} Each venue's depth at those prices, position limits, and on the \omterm{def:m3:blockchains-for-traders:chain}{blockchain} venue the depth of its book or pool. \textbf{9.} Not on a venue that
excludes US persons; the firm's access and legal advice decide. \textbf{10.} That capital does not move freely between them: different users, access rules
and fees keep separate prices, as in the \omterm{def:m3:spot-markets:kimchi}{kimchi premium} (\cref{ch:m3:spot-markets}). \textbf{11.} 16.7\%. \textbf{12.} $0.65/0.58 - 1 = 12.1\%$. \textbf{13.}
Beliefs are uncertain, and a full Kelly stake on an overestimated edge loses heavily; fractions of Kelly trade growth for safety. \textbf{14.} Backtest the
model on past elections against market prices, and compare its calibration and scoring with the market's. \textbf{15.} A 58\% favourite is on the side where
prices tend to be closer to the truth; a longshot view would deserve more suspicion. \textbf{16.} Both: most traders speculate, some hedge exposures to policy
outcomes; the statute's gaming test is a legal line, not an economic one. \textbf{17.} That the CFTC's order prohibiting Kalshi's congressional control
contracts exceeded its authority, because they involve neither unlawful activity nor gaming; it did not decide whether they serve the public interest.
\textbf{18.} The one whose prices are better calibrated over many events; with one election, no one can know from the outcome alone. \textbf{19.}
\emph{Named result:} a gap above about 1.86 cents survives the fees and a month of capital (the 4-cent gap earns 2.14 cents a pair); the Kelly stake on YES at
0.58 for a belief of 65\% is 16.7\% of wealth. \textbf{20.} The market's price of a dollar paid if the event happens: a probability, net of margins, biases and
the cost of capital.
\end{solution}

\begin{solution}{iq:m3:prediction-and-betting-markets:1}
Take reciprocals of decimal odds, which sum to more than one by the \omterm{def:m3:prediction-and-betting-markets:book}{overround}, and remove the margin: proportionally, additively, by a power, or by Shin's
insider model, which put more of the margin on longshots.
\iqlookfor{the \omterm{def:m3:prediction-and-betting-markets:book}{overround} and a model of where the margin sits.}
\end{solution}

\begin{solution}{iq:m3:prediction-and-betting-markets:2}
Longshots return less per unit staked than favourites. Explanations: bettors misperceive small probabilities (prospect theory), risk-loving bettors, and
\omterm{def:m3:prediction-and-betting-markets:book}{bookmakers} shading longshots against insiders (Shin).
\iqlookfor{the fact and at least two explanations.}
\end{solution}

\begin{solution}{iq:m3:prediction-and-betting-markets:3}
Start from a fair probability from models and other venues; quote around it with a spread covering adverse selection and the cost of locked capital; skew
with inventory; hedge on other venues; widen before news and in illiquid hours; cap exposure given the jump to 0 or 1 at resolution.
\iqlookfor{fair value, adverse selection, inventory and the terminal jump.}
\end{solution}

\begin{solution}{iq:m3:prediction-and-betting-markets:4}
Different \omterm{def:m3:prediction-and-betting-markets:event}{resolution sources} and rules, voiding, delayed resolution, venue credit and \omterm{def:m3:blockchains-for-traders:stable}{stablecoin} risk, access restrictions, and capital locked until
resolution.
\iqlookfor{resolution risk first.}
\end{solution}

\begin{solution}{iq:m3:prediction-and-betting-markets:5}
Staking a fraction $f$, wealth becomes $1 + f(o - 1)$ with probability $p$ and $1 - f$ otherwise; maximising $p\ln(1 + f(o - 1)) + (1 - p)\ln(1 - f)$ gives
$f = (po - 1)/(o - 1)$.
\iqlookfor{the log-utility derivation.}
\end{solution}

\begin{solution}{iq:m3:prediction-and-betting-markets:6}
Name primary sources in advance, require resolution after a delay, use bonded proposals with a challenge period and escalation to a broad, economically
bonded vote, publish rules for ambiguous cases, and make the cost of corrupting the vote exceed the value at stake.
\iqlookfor{bonds, delays and escalation, and the cost of attack.}
\end{solution}
